\section{Prospective analysis and resource contract}
\label{app:statistics}

This appendix specifies proposed analysis choices. No confirmation roster,
registered experiment, trained synthesis checkpoint or prospective result is
claimed to exist. Before execution, the campaign must freeze the dataset rights
and family ledger, every component identity, primary comparator, endpoints,
resource envelope, seed schedule, novelty rules, multiplicity and analysis code.
Any change informed by confirmation outcomes creates an exploratory analysis or
requires a new independent lock.

\subsection{Family-weighted estimand and power}

For $F$ independent families with $T_f$ tasks and $R_{ft}$ registered outer
contexts, define the primary difference against locked comparator $B$ by
\begin{equation}
 \Delta_f=\frac{1}{T_f}\sum_{t=1}^{T_f}
   \frac{1}{R_{ft}}\sum_{r=1}^{R_{ft}}
       (u^{\mathrm{IZ}}_{ftr}-u^B_{ftr}),\qquad
 \Delta=\frac{1}{F}\sum_{f=1}^{F}\Delta_f.
\end{equation}
Equal family weighting prevents a parent table with many targets or cuts from
dominating the conclusion. Report task-weighted sensitivity separately. The
primary checkpoint is selected on validation and frozen; replication checkpoints
are not alternatives from which to choose the best confirmation result.

Use 240 families as a planning floor, not a power certificate. With
development-only standard deviation $\sigma$ of paired family differences, a
normal approximation for a true effect $\delta=0.05$, superiority margin
$m=0.02$, two-sided 95\% intervals and 80\% power is
\begin{equation}
 F_{\mathrm{plan}}=\max\!\left\{240,
 \left\lceil\frac{(1.96+0.84)^2\sigma^2}{(0.05-0.02)^2}\right\rceil\right\}.
\end{equation}
For illustration only, $\sigma=0.15$ gives 196 before the floor, whereas
$\sigma=0.25$ gives 545. These are assumptions, not pilot measurements. Cluster
simulations must preserve unequal task counts, failure probabilities and paired
bounded utility; increase the roster if needed. If a feasible budget cannot
support adequate power, declare the study exploratory before opening outcomes.
Do not reduce the margin or stop acquisition after a favorable result.

The error floor is derived from outer-training targets: use
$s=1.4826\,\mathrm{MAD}(y)$ if positive, otherwise
$\mathrm{IQR}(y)/1.349$ if positive, otherwise population standard deviation;
set $\epsilon=\max\{10^{-12}\max(1,|\bar y|),10^{-6}s\}$.
Keep this confirmation utility distinct from historical inner-fold training
labels and their wider clipping range. Publish raw errors and sensitivities to
the clipping and failure floor.

\subsection{Uncertainty and multiplicity}

Use 10,000 fixed-seed paired family bootstrap resamples, preserving every task,
outer context and method within a sampled family. Report the interval method
and resampling seed. These intervals condition on the selected checkpoint;
training-seed effects and dispersion are separate replication evidence. For rare
failures or zero observed failures, supplement percentile intervals with a
conservative family-level bound or justified cluster model. Zero events are not
proof of perfect reliability.

The first confirmatory test is $\Delta>0.02$ against $B$. A fixed-sequence
secondary test compares outcome training with an equal-support static/imitation
model only if the primary gate passes. Practical-baseline pairwise superiority
claims use one prespecified Holm family including unfavorable comparisons.
Reliability and useful-structure rates are separate qualifications, with
simultaneous uncertainty when making multiple standalone claims. History/search
curves and subgroup analyses are secondary unless separately registered and
power-budgeted. Publish every result regardless of gate outcome.

Elo uses declared family-weighted win/tie/loss outcomes and a Bradley--Terry
model on the Elo scale. Freeze opponent set, centering, regularization and
practical tie tolerance. Refit within each family resample. Publish pairwise
matrices, coverage and error magnitudes; absolute ratings from different panels
are not comparable. Noninferiority requires its own paired estimand and margin,
not overlapping Elo intervals.

\subsection{Useful structure and mechanism confirmation}

Let $V$ denote a valid, resource-qualified output. Let $N$ denote out-of-archive
reusable structure after surface normalization and blinded review against every
relevant complete training procedure. Two reviewers assess neutral copies before
seeing outcomes, with adjudication for disagreements; unresolved cases count as
$N=0$ in the conservative primary rate. Preserve original source hashes and
redaction logs. Style can compromise blinding, so record reviewer identity
guesses and uncertainty.

One proposed usefulness rule $Q$ requires a valid ridge receipt, error at least
5\% below fixed ridge and, when $B$ succeeds, error within 5\% of $B$, under
common resource limits. Use the training-derived error floor in these ratios.
If ridge lacks a valid receipt, set $Q=0$ conservatively; if only $B$ fails, retain
the ridge and validity requirements. This is an observed-usefulness predicate,
not a per-program significance test. Report family-weighted
$\mathbf{1}\{V\land N\land Q\}$ over \emph{all} attempts and its marginals.
Proposed sensitivity analyses vary the required ridge improvement over
0/5/10\% while holding the tolerance above $B$ at 5\%, then vary that tolerance
over 0/5/10\% while holding the ridge improvement at 5\%. The primary rule
remains the fixed 5\%/5\% pair. The proposed completion gate requires a lower
95\% family-level bound of at least 95\% on the all-attempt valid-completion
rate. The separate useful-structure gate requires a lower 95\% bound above
10\% on the family-weighted $V\land N\land Q$ rate. Both remain distinct from
the primary utility test.

Choose the lexicographically lowest committed source hash among qualifying
positive-utility cases for mechanism analysis, plus registered parity and failure
examples. Freeze the intervention and matched-complexity control before opening
a separate mechanism lock, initially planned at 60 independent families and
subject to development-only power assessment. A Move42 case requires a positive
paired advantage over the intervention with a lower 95\% bound above zero,
source/provenance qualification and a direction not confined to one acquisition
lineage. The result establishes a bounded mechanism contribution, not universal
originality or a probability of invention.

\subsection{Budgets and event reconciliation}

The reference planning envelope is one CPU-equivalent and 8~GiB for
preprocessing, compilation, fit and prediction; respective elapsed limits are
120, 120, 600 and 60 seconds. Prediction's limit covers all batches together.
Encoding uses a disclosed CPU-only lane or one GPU of at most 24~GiB, capped at
120 seconds. Proposal/lowering/commitment has a 600-second ceiling and one
trajectory. These are proposed common limits, not historical run conditions;
machine/image identities and output limits must be qualified before registration.
Provider-hosted generation is a separate disclosed service lane.

For an i.i.d. reference design, three split seeds and three outer folds produce
nine contexts per task. At 240 families and one task per family this is 2,160
contexts per arm. Eight arms imply 17,280 fits and a maximum fitting allocation
of 2,880 CPU-hours at the 600-second ceiling, before generation, compilation or
search. This is a budget bound, not a runtime forecast. Temporal/grouped families
use appropriate prespecified evaluations rather than forced random folds.

Every scheduled attempt must reconcile to a commitment or typed failure.
Resource caps are checked against actual child-process readbacks. An independent
scorer records prediction hashes and outcomes without returning feedback to the
primary proposer. Infrastructure retries preserve source and retain their event
history. No fallback, missing receipt or dropped difficult task may disappear
from coverage and failure accounting.
